\section*{Conclusion}

L'ensemble des manipulations a permis d'accomplir les trois objectifs fixés pour ce travail pratique :
\begin{itemize}
    \item \textbf{Mise en évidence et caractérisation du pouvoir séparateur de l'œil}, 
    mesuré expérimentalement à environ $\epsilon \approx 10^4$ rad.
    \item \textbf{Étude des grandes caractéristiques du microscope}, 
    notamment par la détermination de son grossissement commercial ($G_c = 100$)
    et de sa puissance intrinsèque ($400 \delta$).
    \item \textbf{Utilisation du microscope dans une situation concrète appliquée à la biologie}, 
    à travers l'étalonnage de l'oculaire micrométrique et la mesure d'échantillons.
\end{itemize}

En réponse à la problématique, 
la compréhension des aspects théoriques du microscope nous a permis de démontrer 
comment ce système optique composé permet de s'affranchir de la limite de résolution naturelle de l'œil. 
En augmentant la puissance optique, 
l'instrument transforme un objet de quelques micromètres 
en une image virtuelle dont le diamètre apparent est supérieur au pouvoir séparateur de l'observateur, 
rendant ainsi possible sa caractérisation quantitative.

La synthèse de nos résultats confirme la validité du modèle théorique, 
bien que l'application biologique ait révélé une anomalie : 
les dimensions relevées ($40$ à $80$ $\mu m$) ont permis d'infirmer l'identification initiale de larve \textit{Pluteus} 
au profit d'un ovocyte d'oursin. 
Cette observation souligne l'importance de la rigueur métrologique 
face aux limites de précision du matériel et aux incertitudes de mise au point.